\hypertarget{adversarial-review-record}{%
\section{Adversarial review record}\label{adversarial-review-record}}

\noindent\textbf{Note on Version 3.} This version adds the reviewer-certified resolution of the first of the
three open programs stated in Version 2 (Section 6.1): at the graph level, the thirteen residual min-cut carriers are
resolved within the declared graph relation (exact same-graph fibers on all thirteen; endpoint orbits disjoint after
complete queue-exhaustive closure under the five declared move classes, modulo terminal-label-fixed exact weighted
isomorphism; local star-mesh moves beyond \(\Delta\)--Y are excluded by cited literature); undeclared graph
transformations remain open. At the state level, six exact finite connected \(C2\_L1\) pairs on those fibers share
identical complete cut fingerprints and identical contracted boundary states while remaining inequivalent under the
declared capacity convention and frozen family relation, separated by an exact generator-proved invariant. This is not
a bulk-geometry theorem. The other two programs (dynamical record formation; persistent invariant candidates) remain
open, with a reviewer-approved research proposal recorded in
\nolinkurl{PROPOSAL_DYNAMICAL_RECORD_STABILITY.md}. Construction and review records:
\nolinkurl{open_programs/} (twelve additional adversarial review rounds; every load-bearing claim independently
recomputed by reviewer and manager). No Version-2 claim is weakened by this version; the abstract, Sections 1, 3, 6,
8--11, and Appendices A, B, D, and E are updated at the corresponding claim and support sites.

\medskip

\noindent\textbf{Note on Version 2.} This version follows a systematic post-publication verification comprising 27
adversarial review rounds and 20 repair or probe constructions, with machine-verifiable records under
\nolinkurl{review_2026/} in the repository. The review retracts the published route-mismatch non-commutation claim, the
forcing claims behind all three predictions, the \(N_{\rm gen}\)-blindness theorem, the universal Born--area
``one-ledger'' claim, the product-parent \(3/23\) control, and the linearized-Einstein leg; it also replaces the apparent
\(\Lambda\) background/\(\kappa\) boundary with the finding that the boundary was a fixed-iteration solver artifact.
The unsound \(11{,}990\)-row gauge carrier is replaced by convention-declared counts of \(1{,}066\) labelled,
\(419\) declared-orbit, or \(195\) primitive-normalized genuinely-chiral structures, and clean separation is restated
as selection of a finite \((\text{structure},\text{breaking orientation})\) branch class rather than a bare structure.
Several results are strengthened by explicit repair, including the branch-complete selection, exact \(SU(5)\) generator
construction, unique content orbit under the declared quotient, materialized mass-rank probe, conditional field-to-mode
provenance, and corrected min-cut LP duality. The complete claim-by-claim mapping from Version 1 to the certified
disposition for Version 2 is \nolinkurl{review_2026/CLAIMS_MAP.md}.

The review history has two distinct phases. Neither is journal peer review or independent replication.

\subsection{Earlier external missing-layer gate}

Version 1 described the external gate as settled. That was incorrect. The surviving record is:

\begin{longtable}[]{@{}p{0.15\linewidth}p{0.27\linewidth}p{0.48\linewidth}@{}}
\toprule
round & recorded disposition & consequence\\
\midrule
\endhead
v7 & approve with changes & approval was conditional; specified changes remained\\
v8 & still needs work & the gate explicitly remained open\\
v9 & unanswered/blank & no final approval was issued\\
\bottomrule
\end{longtable}

Accordingly, the missing-layer atlas is unfit as independent preregistration or blinding evidence. Its cards preserve a
historical research record, but their grades must not be used to override the certified claims map.

\subsection{The 2026 post-publication campaign}

The present correction follows 27 adversarial review rounds and 20 repair or probe constructions completed on
2026-08-26. A manager reproduced review findings; a separate read-only reviewer assessed repairs; implementation was
kept commit-gated. The campaign inspected every load-bearing paper claim, then required repairs to construct missing
objects rather than merely relabel outcomes. Artifacts under \nolinkurl{review_2026/} contain literal dependency pins,
schemas, tables, findings notes, and validators where feasible.

The campaign produced favorable strengthened results as well as retractions. Examples of strengthened constructions are
the conjugation-consistent branch carrier, exact \(SU(5)\) generator action, unique content orbit, real generic mass
ranks, field-to-mode provenance, and explicit min-cut LP duality. Examples of decisive negative findings are the commuting
published route pair, quotient-gauge P1 kernel, token-sensitive record implication, failed honest linear response,
non-shared Born/area composition, and solver-artifact F50 boundary.

The final authority for Version 2 is \nolinkurl{review_2026/CLAIMS_MAP.md}, revision 3, marked CERTIFIED after rounds 26
and 27. It distinguishes KEEP, RESTATE, RETRACT, and OPEN-PROGRAM claims and names every load-bearing condition. This
paper applies that map at each claim site. The campaign is a completed adversarial pass for this release; it is not a
claim that the framework or paper has been independently validated.

\subsection{Audit interpretation}

The honest methodological conclusion is mixed. The repository contains a real self-correction trail and several
machine-verifiable repairs. It also contains historical validators that regenerate in place, validate stored summaries,
or used a tautological pin. Passing the repository's validators is therefore necessary evidence of consistency, not a
tamper-evident freeze or a substitute for semantic review. Appendix B states the operational caveats for reruns.
